% ==============================================================
% PAQUETES Y CONFIGURACIÓN BASE
% ==============================================================
\usepackage[utf8]{inputenc}
\usepackage[T1]{fontenc}
\usepackage[x11names, dvipsnames, svgnames]{xcolor} % Carga indispensable de colores

% --- Fuente Palatino ---
\usepackage{mathpazo}
\linespread{1.05}
\usepackage{textcomp}

% --- Matemáticas ---
\usepackage{amsmath, amssymb, amsthm, mathtools, cancel}

% --- Diseño, Tablas y Flotantes ---
\usepackage{enumitem, booktabs, fancyhdr, titlesec, tabularx}
\usepackage{graphicx, float, subcaption}
\usepackage{stfloats} % Evita que las imágenes de 2 columnas floten sin control

% --- Cajas y Diagramas ---
\usepackage[most]{tcolorbox}
\usepackage{tikz, tikz-cd}

% --- Referencias cruzadas ---
\usepackage{hyperref}
\hypersetup{
    colorlinks   = true,
    linkcolor    = MidnightBlue,
    citecolor    = Teal,
    urlcolor     = DarkSlateGray
}
\usepackage[spanish]{cleveref}

% ==============================================================
% PALETA DE COLORES
% ==============================================================
\colorlet{colordef}{MidnightBlue}
\colorlet{colorteor}{Teal}
\colorlet{colorej}{DarkSlateGray}
\colorlet{colorobs}{DimGray}

% ==============================================================
% ESTILOS DE CAJAS (tcolorbox)
% ==============================================================
\tcbset{
    theorembase/.style={
        fonttitle            = \bfseries\sffamily\small,
        separator sign       = {\ ---},
        description delimiters = {(}{)},
        terminator sign      = {.},
        before upper         = {\parindent 1.5em},
        breakable,
        enhanced,
        attach boxed title to top left = {yshift = -2mm, xshift = 4mm},
        boxed title style    = {sharp corners, size = small},
        sharp corners,
        top = 4mm,
    },
    defstyle/.style={
        theorembase,
        colback      = colordef!4!white,
        colframe     = colordef!50!black,
        colbacktitle = colordef!70!black,
        coltitle     = white,
    },
    teorstyle/.style={
        theorembase,
        colback      = colorteor!4!white,
        colframe     = colorteor!55!black,
        colbacktitle = colorteor!75!black,
        coltitle     = white,
    },
    ejstyle/.style={
        theorembase,
        colback      = colorej!4!white,
        colframe     = colorej!45!black,
        colbacktitle = colorej!65!black,
        coltitle     = white,
    },
    obsstyle/.style={
        enhanced,
        breakable,
        fonttitle       = \bfseries\sffamily\small,
        separator sign  = {\ ---},
        description delimiters = {(}{)},
        terminator sign = {.},
        before upper    = {\parindent 1.5em},
        attach boxed title to top left = {yshift = -2mm, xshift = 4mm},
        boxed title style = {sharp corners, size = small,
                             colback = colorobs!60!black, coltitle = white},
        colback  = colorobs!4!white,
        colframe = colorobs!4!white,
        borderline west = {2.5pt}{0pt}{colorobs!70!black},
        top = 4mm,
    },
}

% ==============================================================
% AMBIENTES DE TEOREMAS
% ==============================================================
\newtcbtheorem[number within = section]{teorema}{Teorema}{teorstyle}{thm}
\newtcbtheorem[use counter from = teorema]{lema}{Lema}{teorstyle}{lem}
\newtcbtheorem[use counter from = teorema]{proposicion}{Proposición}{teorstyle}{prop}
\newtcbtheorem[use counter from = teorema]{corolario}{Corolario}{teorstyle}{cor}
\newtcbtheorem[number within = section]{definicion}{Definición}{defstyle}{def}
\newtcbtheorem[number within = section]{ejemplo}{Ejemplo}{ejstyle}{ej}
\newtcbtheorem[number within = section]{observacion}{Observación}{obsstyle}{obs}
\newtcbtheorem[number within = section]{ejercicio}{Ejercicio}{obsstyle}{exer}

\renewcommand{\proofname}{Demostración}

% ==============================================================
% MACROS MATEMÁTICAS
% ==============================================================
\newcommand{\N}{\ensuremath{\mathbb{N}}}
\newcommand{\Z}{\ensuremath{\mathbb{Z}}}
\newcommand{\Q}{\ensuremath{\mathbb{Q}}}
\newcommand{\R}{\ensuremath{\mathbb{R}}}
\newcommand{\C}{\ensuremath{\mathbb{C}}}

\newcommand{\eps}{\varepsilon}
\newcommand{\del}{\delta}

\DeclarePairedDelimiter{\abs}{\lvert}{\rvert}
\DeclarePairedDelimiter{\norm}{\lVert}{\rVert}
\DeclarePairedDelimiter{\inner}{\langle}{\rangle}
\DeclarePairedDelimiter{\ceil}{\lceil}{\rceil}
\DeclarePairedDelimiter{\floor}{\lfloor}{\rfloor}

\DeclareMathOperator{\sen}{sen}
\DeclareMathOperator{\tg}{tg}
\DeclareMathOperator{\Img}{Im}
\DeclareMathOperator{\Ker}{Ker}
\DeclareMathOperator{\rango}{rango}

% ==============================================================
% FORMATO DE SECCIONES (titlesec)
% ==============================================================
\titleformat{\section}
    {\large\bfseries\sffamily}
    {\thesection.}{0.6em}{}
    [\vspace{2pt}\rule{\linewidth}{0.4pt}\vspace{-4pt}]

\titleformat{\subsection}
    {\normalsize\bfseries\sffamily}
    {\thesubsection.}{0.5em}{}

\titleformat{\subsubsection}
    {\normalsize\itshape}
    {\thesubsubsection.}{0.5em}{}

\titlespacing{\section}{0pt}{14pt plus 2pt minus 1pt}{6pt}
\titlespacing{\subsection}{0pt}{10pt plus 1pt minus 1pt}{4pt}

% ==============================================================
% ENCABEZADO Y PIE DE PÁGINA (CORREGIDO PARA EVITAR SOBREESCRITURA)
% ==============================================================
\pagestyle{fancy}
\fancyhf{}
\fancyhead[L]{\small\scshape MA2003B Análisis Multivariado}
\fancyhead[R]{\small\scshape Predicción de $PM_{10}$ en la ZMM}
\fancyfoot[C]{\small\thepage}

\setlength{\headheight}{14pt}
\renewcommand{\headrulewidth}{0.4pt}
\renewcommand{\footrulewidth}{0pt}

\fancypagestyle{plain}{%
    \fancyhf{}
    \fancyfoot[C]{\small\thepage}
    \renewcommand{\headrulewidth}{0pt}
}

\setlength{\parskip}{4pt}
\setlength{\parindent}{1.2em}

% Evitar que los títulos queden huérfanos al pie de columna
\usepackage[bottom]{footmisc}
\usepackage{needspace}

% Restricción severa para que los flotantes no se acumulen al final
% (extraplaceins ya no está en TeX Live; se usa placeins con \FloatBarrier
% automático antes de cada \section)
\usepackage{placeins}
% Flotantes de la sección actual antes de pasar a la siguiente
\let\Section\section
\renewcommand\section[1]{\FloatBarrier\Section{#1}}
\maxdeadcycles=200

% ==============================================================
% CONTROL DE FLOTANTES EN DOS COLUMNAS (REPARACIÓN DE MAQUETACIÓN)
% ==============================================================
% Permite más espacio para texto y menos bloqueo por figuras
\renewcommand{\topfraction}{0.85}
\renewcommand{\bottomfraction}{0.7}
\renewcommand{\textfraction}{0.15}
\renewcommand{\floatpagefraction}{0.66}
\renewcommand{\dbltopfraction}{0.66}

% Fuerza a que los encabezados mantengan su posición sin ser empujados por flotantes
\titleformat{\section}[block]{\large\bfseries\sffamily}{\thesection.}{0.6em}{}[\vspace{2pt}\rule{\linewidth}{0.4pt}\vspace{-4pt}]